\documentclass[11pt]{article}

\usepackage[a4paper,margin=2.5cm]{geometry}
\usepackage{amsmath,amssymb}

\title{SMC per un masked diffusion language model\\
\large Nota matematica minimale}
\author{}
\date{}

\begin{document}
\maketitle

\section{Modello base}

La generazione del masked diffusion model e' una catena di Markov
\[
X_0=[\mathrm{MASK}]^L,\qquad X_t\sim K_t(\,\cdot\mid X_{t-1}),
\quad t=1,\ldots,T.
\]
La probabilita' congiunta di una traiettoria e'
\[
P_0(x_{0:T})
=\prod_{t=1}^{T}K_t(x_t\mid x_{t-1}).
\]
In un singolo passo, i token ancora mascherati possono essere campionati
condizionatamente in parallelo. Questo non rende indipendenti i token della
completion: ogni distribuzione dipende dall'intero stato precedente. La
probabilita' terminale e'
\[
p_0(x_T)=\sum_{x_1,\ldots,x_{T-1}}P_0(x_{0:T}),
\]
ma l'algoritmo non deve calcolare esplicitamente questa somma.

\section{Distribuzione tilted}

Sia $R(x_T)$ il reward terminale e sia $\lambda\geq 0$. Vogliamo campionare
dalla distribuzione
\[
\pi_\lambda(x_T)
=\frac{p_0(x_T)e^{\lambda R(x_T)}}{Z_\lambda},
\qquad
Z_\lambda=\mathbb E_{p_0}[e^{\lambda R(X_T)}].
\]

Scegliamo checkpoint $0=t_0<t_1<\cdots<t_m=T$. Per $k<m$, completiamo
deterministicamente lo stato parziale $X_{t_k}$ tramite argmax e calcoliamo
un reward provvisorio $R_k(X_{t_k})$. Al terminale usiamo la completion
realmente campionata:
\[
R_m(X_T)=R(X_T).
\]
Definiamo il log-potenziale
\[
\Phi_k(X_{t_k})=\beta_k R_k(X_{t_k}),
\qquad \beta_0=0,\quad \beta_m=\lambda.
\]

\section{Peso SMC e resampling}

Ogni particella viene propagata con lo stesso kernel $K_t$ del modello base.
Al checkpoint $k$, il suo peso viene moltiplicato per
\[
G_k
=\exp\!\left(
\Phi_k(X_{t_k})-\Phi_{k-1}(X_{t_{k-1}})
\right).
\]
Equivalentemente,
\[
\log w_k
=\log w_{k-1}
+\Phi_k(X_{t_k})-\Phi_{k-1}(X_{t_{k-1}}),
\]
seguito dalla normalizzazione dei pesi.

Se l'ESS
\[
\operatorname{ESS}=\frac{1}{\sum_i (W_k^{(i)})^2}
\]
scende sotto la soglia, si estraggono i genitori con probabilita'
$W_k^{(i)}$ e si riportano i pesi a $1/N$. L'informazione dei pesi precedenti
rimane nella molteplicita' dei discendenti.

Il prodotto dei potenziali lungo ogni genealogia telescopa:
\[
\prod_{k=1}^{m}G_k
=\exp\!\left(\Phi_m(X_T)-\Phi_0(X_0)\right)
=e^{\lambda R(X_T)}.
\]
Di conseguenza, i reward provvisori influenzano il resampling e quindi
l'efficienza, ma non modificano il target terminale.

\section{Perche' non serve la joint probability della completion}

In generale, un peso di importance sampling contiene il rapporto tra il
kernel del target $K_t$ e il kernel di proposta $Q_t$. Qui la propagazione
non viene modificata:
\[
Q_t=K_t
\qquad\Longrightarrow\qquad
\frac{K_t}{Q_t}=1.
\]
Le probabilita' del masked diffusion model sono quindi gia' incluse nel
campionamento della traiettoria e si cancellano nel rapporto. Rimane soltanto
il rapporto tra i potenziali $G_k$.

Questa giustificazione richiede che la propagazione resti esattamente quella
del modello base. Se si modificassero logits, temperatura o kernel di
demasking, sarebbe necessario includere il corrispondente rapporto
$K_t/Q_t$ nei pesi.

\section{Stima dell'evento raro}

Sia $E$ l'evento raro. Dalla distribuzione tilted vale l'identita'
\[
p_0(E)
=\frac{
\mathbb E_{\pi_\lambda}
\left[e^{-\lambda R(X)}\mathbf 1_E(X)\right]
}{
\mathbb E_{\pi_\lambda}
\left[e^{-\lambda R(X)}\right]
}.
\]
Se $(X_i,W_i)$ e' la popolazione SMC terminale, definiamo
\[
\widetilde W_i
=\frac{W_i e^{-\lambda R(X_i)}}
{\sum_{r=1}^{N}W_r e^{-\lambda R(X_r)}}
\]
e stimiamo
\[
\widehat p_\lambda(E)
=\sum_{i=1}^{N}\widetilde W_i\mathbf 1_E(X_i).
\]
Con $J$ popolazioni indipendenti e pesi uniformi,
\[
\widehat p(E)=\frac{1}{J}\sum_{j=1}^{J}\widehat p_{\lambda_j}(E).
\]

\section*{Conclusione}

L'SMC e' corretto perche' opera sulla probabilita' congiunta della
\emph{traiettoria}, usa il kernel originale come proposta e applica
potenziali che telescopano fino a $e^{\lambda R(X_T)}$. Non calcola la
probabilita' puntuale $p_0(x_T)$ di una completion, e non ne ha bisogno.
Per un numero finito di particelle, lo stimatore self-normalized non e'
esattamente unbiased; e' consistente e la sua affidabilita' dipende dalla
non degenerazione dei pesi corretti.

\end{document}
